%!TEX program = lualatex
\documentclass[ngerman,headheight=70pt]{scrartcl}
\usepackage[ngerman]{babel}
\usepackage{blindtext}
\usepackage[utf8]{luainputenc}
\usepackage{scrlayer-scrpage}
\usepackage[T1]{fontenc}
\usepackage{fontspec}
\usepackage{multicol}
\usepackage{enumitem}
\usepackage[german=quotes]{csquotes}
\usepackage[bottom=2cm,footskip=8mm]{geometry}
\setmainfont{Source Sans Pro}
\MakeOuterQuote{"}

\setlength{\parindent}{0pt}
\setlength{\parskip}{6pt}

\pagestyle{plain.scrheadings}
\thispagestyle{scrheadings}
\lohead{\vspace*{-1.5cm}\includegraphics[scale=0.4]{../images/stupa_logo_rot.pdf}}
\rohead{\\Präsidium des\\Studierendenparlaments}

\lofoot[Universität Hamburg · Präsidium des Studierendenparlaments\\\hrule Von-Melle-Park 5 · D-20146 Hamburg]{Universität Hamburg · Präsidium des Studierendenparlaments\\\hrule Von-Melle-Park 5 · D-20146 Hamburg}
\cofoot[]{}
\cohead[\thepage]{}

\renewcommand{\thesubsection}{TOP \arabic{subsection}}
\renewcommand{\thesubsubsection}{\arabic{subsubsection}.}
\setcounter{subsection}{-1}

\begin{document}
    UHH · StuPa-Präsidium · Von-Melle-Park 5 · D-20146 Hamburg

    %%% PAGE ONE (EDIT LATER) %%%%
    % Hier den Titel des Protokolls eintragen
    \section*{Protokoll der 2. Sitzung des Studierendenparlaments vom 12. Mai 2016}

    % Protokollanten eintragen
    \textbf{Protokoll: Gunhild Berdal}\\
    \textbf{Ort: S08 in der ehemaligen HWP}\\
    \textbf{Beginn: 18.40 Uhr}\\
    \textbf{Ende: 0.02 Uhr}

    % Wohl die komplizierteste Stelle im Protokoll
    % Als Faustregel sollten nicht mehr als 3 Namen auf einer Zeile stehen
    % Einfach an dieser Vorlage orientieren
    % Im Idealfall ist hier nicht viel zu ändern, da kaum jemand fehlt.
    % Rücktritte müssen natürlich pro Protokoll angepasst werden
    \vspace{0.5cm}
    \begin{tabular}{ll}
        Anwesend: & \\
            RCDS (5 Sitze): & Ramon Weilinger, Antonia Niecke, Ramin Shakiba, \\
                            & Jennifer Maack, Benjamin Welling \\
             CampusGrün (14 Sitze): & Geoffrey Youett, Philipp Droll, Melf Johannsen, \\
                                   & Tahnee Herzig, Mario Moldenhauer, Jim Martens,\\
                                   & Mirzo Ulugbek Khatamov, Armin Günther, Martin Sievert\\
             Bier-Liste (2 Sitze) : & Jakob Pape \\
             WiWi (2 Sitze): & Claas-Friso Hente \\
             Unicorns (5 Sitze): & Johannes Peplow, Annkathrin Löffler, \\
                                 & Andreas Hartkamp, Marielle Hermstrüwer \\
             Liste LINKS (3 Sitze): & Gunhild Berdal, Till Petersen, Sinah Mielich \\
             HWP (2 Sitze): & Samet Gunay, Ajdina Karahasan \\
             MIN (4 Sitze): & \\
             SDS* (3 Sitze): & Mena Winkler, Jacob Petersein, Artur Brückmann \\
             Bart-LISTE (2 Sitze): & Dominic Laumer \\
             LHG (1 Sitz): & Tobias Heisig \\
             harte zeiten (1 Sitz): & Tobias Berking \\
             Jura (1 Sitz): & \\
             AL (2 Sitze): & Karima Schulze, Henri Weber \\
            & \\
        Entschuldigt: & Freya Schmitz, Yasemin Günther, Svenja Horn,\\
                      & Elena Rysikova (alle Campus Grün), Lotte Rullkötter, Ailina Salten,\\
                      & Thea Wahlers (alle MIN), Katharina Kucza (Unicorns), Timo Zeimet (Bart) \\
                                &\\
        Unentschuldigt abwesend: & Johann Baumhoefener (Jura), Laura Franzen (CG), Karen Martirosian (Bier), \\
                                & Kay Zöllmer (WiWi), Jan Detampel (MIN) \\
                                &\\
        Rücktritte: & \\
    \end{tabular}
    \newpage
    %%% Vorgeschlagene Tagesordnung %%%
    \underline{Vorgeschlagene Tagesordnung}
    \begin{enumerate}[label={\textbf{Top \theenumi}},leftmargin=*]
        % folgende Zeile bei Sitzungsfortsetzungen einkommentieren und x durch
        % die Zahl des zuletzt behandelten TOPs ersetzen
        %\setcounter{enumi}{x}
        \item Alle Frauen*-Wahl
            \begin{enumerate}
                \item Bestätigung der Wahlniederschrift (V1617-020)
                \item Bestätigung der Referent*innen (V1617-019)
            \end{enumerate}
        \item Gründung des FSR kath. Theologie (V1617-017) (10 Min.)
        \item Antrag auf VV (V1617-018) (30 Min.)
        \item EM 2016 (V1617-014) (30 Min.)
        \item Verlegung der StuPa-Wahl in das Sommersemester (V1617-015) (30 Min.)
        \item Umgang mit der AfD (V1617-016) (30 Min.)
        \item Verschiedenes
    \end{enumerate}

    \newpage

    % Hier beginnt Teil A. Bei Fortsetzungen kommt hier nur der Teil über
    % die Feststellung der endgültigen Fassung der Tagesordnung.
    % Für sonstige Sitzungen ist dies bereits das Gerüst für die Formalia.
    % Man muss nur noch die protokollierten Inhalte eintragen.

    {\Large \underline{Teil A}}

    \subsection{Formalia}

    Jim eröffnet die Sitzung um 18.40 Uhr.

    \subsubsection{Geschäftsbericht Präsidium}

    Öffentlichkeitsarbeit

    Fegebank

    Sitzungsplanung:
    Am 23. Juni findet die Tagung in der Erz.Wiss zu Noten, eine APH- und eine
    Wissenschaftausschusssitzung statt. Der Termin ist daher für eine StuPa-Sitzung
    nicht so gut geeignet. Es wird daher überlegt die Sitzung vom 23. Juni auf den
    16. Juni vorzuverlegen.

    Daraus würden sich diese Sitzungstermine ab dem 16. Juni ergeben:
    \begin{itemize}
        \item 16. Juni
        \item 30. Juni
        \item 14. Juli (bei Bedarf)
    \end{itemize}

    \subsubsection{Anfragen an das Präsidium}

    Keine

    \subsubsection{Geschäftsbericht AStA}

    Philipp: AStA konstituieren, Referatszuschnitte, Schlüssel usw. Viel Raum
    eingenommen.
    Veranstaltung T. Bultmann (BdWi). Steigende Tendenz … Abhängigkeit von
    Drittmitteln. Strukturelle Scheinautonomie der Hochschulen. Es sind dieselben
    Hochschulen, die Mittel einwerden.
    Akkreditierung von Hochschulen und Studiengängen. Noch keine Linie:
    Einheitlichkeit eher positiv, negativ die Ba/MA-Studiengänge

    Franziska: Demo zu Holzwirtschaft: Uni für alle statt Exellenz, viele Plakate
    und aktive Studierende. Demo gegen die Exellenzinitiative am 14. Juni und VV
    Gebührensatzungen von Master of International Taxation (Master of Panama Papers),
    Master of legal Studies (Troika-Master) und Master der kriminologischen Studis
    Rahmenprüfungsordnung: Diskussionstag in der vorlesungsfreien Zeit
    Studierende im AS: Ausschuß für internationale Angelegenheiten

    Timo: Erste Überlegungen für den Haushalt

    Artur: Zitat von wem? "Meine Einschätzung aber .. " Horst Köhler 2010
    zurückgetreten. Bildung mündiger Menschen statt Gehhorsam, Bundeswehr ausgeladen.
    Gesundheit für alle, Veranstaltung mir Mitarbeitern der Charité, die mehr
    Personal für die Pflege durchgesetzt haben
    Exellenzinitiative

    Armin: Öko: kritische Veranstaltung zu Tierversuchen, Recycling in der Stadt.
    Geoffrey: Öff: Viel org. Kram, besser mit den KommilitonInnen in Kontakt kommen

    Kim: EM problematisieren und kritisieren

    Janis: Kulturkurse funktionieren, was nicht so gut läuft, wird geprüft

    Mario: Viele Projekte, keine Konzepte. Seld: Altersfeststellung

    Oliver: Ausschuß gegen Rechts: Lesung "Bücherverbrennung nie wieder" und
    "Tag gegen Rechts". Gutes Wetter, viel Besuch, gute Texte, viel gelacht.
    Nicht nur Pazifismus und Antimilitarismus und Sozialismus und so, sondern auch
    Humor war den Nazis ein Dorn im Auge.

    \subsubsection{Anfragen an den AStA}

    EM gegen Rechts?

    Antworten Antidiss: Kulturbereich gegen Rechts

    Gunhild: Gemeinsame Idee, gemeinsame Vorhaben?

    Philipp: Exellenzinitiative

    \subsubsection{Dringlichkeitsanträge des AStA}

    Keine

    \subsubsection{Aktuelle Stunde (falls entsprechender Antrag vorliegt)}

    Es lag kein Antrag vor.

    \subsubsection{Feststellung der endgültigen Fassung des Teils B der Tagesordnung}

    Till: Top 5 und 6 tauschen; M:9:0

    Claas-Friso: Top 5 als 4; 2:M:3 \\
                 Top 5 und 1 tauschen; 8:m:1

    Geoffrey: Wahl nach vorne ziehen falsch, dafür haben wir Zeit

    Till: TO heute bewältigen: Zuerst formal, dann politischer, dann unsere eigene Formalia

    %%% BESCHLOSSENE TAGESORDNUNG %%%
    \underline{Endgültige Fassung der Tagesordnung}
    \begin{enumerate}[label={\textbf{Top \theenumi}},leftmargin=*]
        % folgende Zeile bei Sitzungsfortsetzungen einkommentieren und x durch
        % die Zahl des zuletzt behandelten TOPs ersetzen
        %\setcounter{enumi}{x}
        \item Alle Frauen*-Wahl
            \begin{enumerate}
                \item Bestätigung der Wahlniederschrift (V1617-020)
                \item Bestätigung der Referent*innen (V1617-019)
            \end{enumerate}
        \item Gründung des FSR kath. Theologie (V1617-017) (10 Min.)
        \item Antrag auf VV (V1617-018) (30 Min.)
        \item EM 2016 (V1617-014) (30 Min.)
        \item Umgang mit der AfD (V1617-016) (30 Min.)
        \item Verlegung der StuPa-Wahl in das Sommersemester (V1617-015) (30 Min.)
        \item Verschiedenes
    \end{enumerate}

    \subsubsection{Feststellung der Beschlussfähigkeit}

    31/47 anwesend

    \subsubsection{Genehmigung der Protokolle der vorangegangenen Sitzungen}

    Till problematisiert die verkürzte Zusammenfassung einzelner Redebeiträge,
    im ersten Protokoll z.B. S. 10 Jacob sagt, es sei kein Problem, in beiden
    Ausschüssen zu sein auf Nachfrage von Seld.
    Wir arbeiten ein. Bestätigung des Protkolls.

    2. Teil: Till beantragt Vertagung – kein Widerspruch

    \textit{Präsidiumspause bis 20.05 Uhr}

    \vspace{1cm}
    {\Large \underline{Teil B}}

    % Teil B beginnt
    % Dieser wird sich für jede Sitzung unterscheiden.
    % Jeder TOP ist eine subsection.
    % Vorhandene Beispiele anfangs entfernen und durch tatsächliche TOPs ersetzen

    % Folgende Zeile bei Sitzungsfortsetzungen einkommentieren und x mit
    % der Zahl des zuletzt behandelten TOPs ersetzen.
    %\setcounter{subsection}{x}
    \subsection{Alle Frauen*-Wahl}

    a) Bestätigung der Wahlniederschrift

    Franziska, Mena und Paula stellen die Wahlniederschrift vor.
    Sie wird mit m:0:2 angenommen.

    b) Bestätigung der Referent*Innen

    Referent*Innenvorschlag:
    \begin{itemize}
        \item Melina Toelle
        \item Louisa Escher
    \end{itemize}

    Melina Toelle (Referentin): allgemein zum Selbstverständnis des AlleFrauen*Referats.
    Vom feministischen Standpunkt: Gesellschaftliche Ungleichheit wird durch Ausbeutung
    und Machtgefälle im Kapitalismus stabilisiert. Patriachale Strukturen aufbrechen.
    Konkret im Studienalltag: Wer macht den Mund auf?

    Louisa Escher (Stellvertreterin): Wie (sozial) offen ist die Uni? Konkurrenz
    verschärft die Ungleicheit. Tag gegen Rechts mitgestalten. Sexulisierte Gewalt
    sichbar machen, aber auch ein solidarisches Miteinander schaffen.

    Karima: Arbeitet Ihr zu leeren Stellen, die nicht mit Frauen besetzt werden?

    Louisa:

    Till: Ausbeutungsverhältnis zwischen Mann und Frau. Konkurrenz gilt allen.
    Positive Diskriminierung verdoppelt die Ungleichheit. Mann gilt als Ausbeuter
    von Frauen. Karte für den Campus: Enttabuisierung richtig, aber man fühlt sich
    auf dem Campus nicht sicher.

    Louisa: Strukturen der Ausbeutung tut niemandem gut.
    Wichtig, daß man nicht alleine ist.

    Melina: Betonung liegt auf den Strukturen.

    Melanie: Patriotismus und Sexismus bei der EM

    Louisa:

    Till: Leistungsideologie im Fußball, Leistung = männlich.
    Beide – Mann und Frau – sind ausgebeutet, internes Hierarchieverhältnis

    Offene Abstimmung: m:0:9 bestätigt

    \textit{Zeitansage: 20.50 Uhr}

    \subsection{Gründung des FSR kath. Theologie}

    Geoffrey stellt den Antrag vor.

    Till: Das katholische Verständnis von Faschschaftsratarbeit nicht hilfreich.
    Jemand vom AStA sollte mit ihnen sprechen, wie man FSR-Arbeit machen kann.
    Warum fällt es in Hamburg mit dem Laizismus so schwer? Nur Leute einer
    Konfession dürfen in den Schulen Religion unterrichten.

    Antonia, Jennifer: FSR erstmal gründen, dann diskutieren

    Änderungsantrag: Nicht der "Fachschaftsrat", sondern die Fachschaft soll eingerichtet werden.

    Mit m:0:3 angenommen

    \subsection{Antrag auf VV}

    Franziska: Kritik an der Exellenzinitiative nimmt bundesweit zu. Ist überall
    Thema, z.B. Diskussion über Studiengang "Neuropsychologie". Mit so vielen wie
    möglich sprechen.

    Ramin: Soll nur verhindert werden, daß die Exini weitergeführt wird, oder
    soll es mehr Geld für die Uni geben?

    Geoffrey: Konzept absurd. Fortsetzung = Verschärfung. Nicht nur Exellenzcluster,
    sondern Exellenzunis.

    Ramin: Gegen Kürzungen. Aber wie kann man gegen die Unterstützung wichtiger
    Forschung, z.B. Klimaforschung, sein?
    Naturwissenschaften brauchen mehr Geld, Desy sehr erfolgreich.

    Till: RCDS möge es kritischer sehen. Exellenz selbstreferenziell.
    Gleichheitsphobie sensationell. Nicht Gleichmacherei, sondern nach Bedarf.
    Kriterium: Quantität und Verwertungstauglichkeit.

    Abstimmung Antrag: m:8:0

    \textit{Präsidiumspause bis 21.50 Uhr}

    \subsection{EM 2016}

    Jakob Pape stellt den Antrag vor: In Israel palästinensiche und israelische
    Mädchen zusammengespielt, sprechen dieselbe Sprache. Hier Flüchtlinge und
    andere. Das Verbindende ist wichtig.
    Medien sind auch mal kritisch, z.B. in Brasilien. Veranstaltung: zuerst die
    Spiele angucken und dann kritisch diskutieren.

    Benjamin:

    Oliver: Bei WM und EM kommt alles zusammen, was an der bürgerlichen
    Gesellschaft scheiße ist: Rassismus, Antisemitismus, "Du bist Deutschland".
    Verdrängung der sozialen Frage, ich kenne keine Parteien mehr, nur noch Deutsche.

    Geoffrey: Gedicht. Das Nationale ist immer falsch

    Till: Großevents auf dem Campus

    Benjamin: Menschen neigen dazu, Gruppen zu bilden. EM verbieten hilft nicht
    gegen Antisemitismus und führt nicht zu Frieden.

    Geoffrey: Es geht nicht ums Verbieten. Was machen solche Großevents mit uns?
    Wohlfühlen, um die soziale Spaltung zu vertuschen. Jetzt genau das Falsche,
    Patriotismus zu befördern.

    Till: Die Menschheit ist eine Gruppe. Massengruppenbildung nicht falsch, z.B.
    viele protestieren gegen Entlassungen.
    Zivilcourage bei Spielen, wo Nationalismus gesungen wird, löst das Problem nicht.

    Melf: Verlängerung der Sitzung bis nach TOP 7. Kein Widerspruch

    Die Frage entsteht, wie mit dem Ersetzungsantrag umzugehen ist: ist es ein
    Änderungsantrag oder wird alternativ abgestimmt?
    Es wird sich darauf geeinigt, so wie mit dem Antrag des RCDS zu AfD auf der
    letzten Sitzung zu verfahren – den Ersetzungsantrag zuerst abstimmen – das
    Präsidium soll die Frage dann grundsätzlich klären.

    Der Ersetzungantrag wird mit m:6:3 angenommen.

    Der Gesamtantrag wird mit m:7:0 angenommen.

    \textit{Präsidiumspause bis 23 Uhr}

    \subsection{Umgang mit der AfD}

    Ramin stellt den Antrag vor

    Till: Gibt es ein ernstgemeintes Anliegen? Jeder Satz falsch... Änderungsantrag für Tag gegen Rechts

    Franziska: Satire und Humor, Erich Kästner

    Till: Die "Sorge um innere Sicherheit"
    Artur

    Jakob Pape: Populismus

    Golnar: Deutschland soll in der internationalen Konkurrenz behaupten

    GO-Antrag Melf auf Schließung der Redeliste: mit 11:9:4 angenommen.

    Till: Ideologie der inneren Sicherheit, wenn alle ordentlich sind, funktioniere
    der Kapitalismus. Dagegen ordnungspolitische Maßnahmen, Polizei und Militär

    Änderungsantrag LL/harte Zeiten/SDS: 'Das StuPa bekräftigt, dass der
    Ausschuss gegen Rechts in Zusammenarbeit mit dem AStA einen "Tag gegen Rechts"
    entwickeln, planen und durchführen mögen und wird dies nach Kräften unterstützen.
    Insbesondere möge vorbereitet werden, dass sich an diesem Tag mit der AfD
    kritisch-analytisch auseinandergesetzt wird. Auf der Grundlage der Erkentnisse
    des "Tages gegen Rechts" soll der Ausschuss gegen Rechts eine politische
    Stellungnahme gegen die AfD ausarbeiten, die dem StuPa zur Beschlussfassung
    vorgelegt wird.'

    Mit m:7:0 angenommen

    Gesamtabstimmung: m: 0: 6 angenommen

    \subsection{Verlegung der StuPa-Wahl in das Sommersemester}

    \textit{GO-Antrag Melf: Vertagung des TOPs auf die nächste Sitzung: m: 5: 3 angenommen}

    \subsection{Verschiedenes}

    Franziska: Neue Veranstaltung zur Zivilklausel am 6.6.16

\end{document}
